\begingroup
\renewcommand{\arraystretch}{1.45}
\SourceLeaf{312}
\Needspace{35\baselineskip}
\InterpolationArticle{25}
If the function $X$ does not stop at the terms $\cos mx$, $\sin mx$ but continues further, and the coefficients of the following terms are still too considerable to be neglected, the coefficients $\gamma,\gamma',\delta'$ etc. in the function $X'$, which satisfies the $\mu$ values of the function $X$ forming a complete period for $x=a,b,c,d$ etc., will differ notably from the corresponding coefficients $\alpha,\alpha',\beta'$ etc. Thus, if instead of the period $x=a,b,c,d$ etc. we treat another period of $\mu$ terms, say for $x=a',b',c',d'$ etc., in the same way, and develop a function of the $m^{\mathrm{th}}$ order satisfying these new values of the function $X$, and represent it, just as we did $X'$, by the expression
\[
\begin{aligned}
&\gamma+\gamma'\cos x+\gamma''\cos2x+.\,.\,.+\gamma^m\cos mx\\[5pt]
&\phantom{\gamma}+\delta'\sin x+\delta''\sin2x+.\,.\,.+\delta^m\sin mx
\end{aligned}
\]
we shall now obtain values for the coefficients $\gamma,\gamma',\delta'$ etc. notably different from those we found before. This shows how these coefficients may be regarded as \textit{variable}. Indeed, since it is clear from Articles 21, 23 that each is determined by the angle $\mu a'$ in the same way as the earlier values were by the angle $\mu a$, each coefficient can be considered a function of the undetermined angle $y$, just as $X$ is a function of the angle $x$. \Needspace{34\baselineskip}
Now let us suppose that $\mu\nu=\pi$ values of the function $X$ forming a complete period are known for
\[
\begin{aligned}
&x=a,a',a'',.\,.\,.b,b',b'',.\,.\,.c,c',c'',.\,.\,.d,d',d''.\,.\,.\text{ etc.}
\end{aligned}
\]
so that
\[
\begin{aligned}
&a'=a+\tfrac1\pi360^0.\\[5pt]
&a''=a+\tfrac2\pi360^0\text{ etc.}\\[5pt]
&b=a+\tfrac\nu\pi360^0=a+\tfrac1\mu360^0,\\[5pt]
&b'=a+\tfrac{\nu+1}\pi360^0=a'+\tfrac1\mu360^0\text{ etc.}\\[5pt]
&c=a+\tfrac{2\nu}\pi360^0=a+\tfrac2\mu360^0\text{ etc.}
\end{aligned}
\]
and that it is evident that this period of $\pi$ terms can be divided into $\nu$ periods of $\mu$ terms for
\[
\begin{aligned}
&x=a,\ b,\ c,\ d\text{ etc.}\\[5pt]
&x=a',\ b',\ c',\ d'\text{ etc.}\\[5pt]
&x=a'',\ b'',\ c'',\ d''\text{ etc.}\\[5pt]
&\text{etc.}
\end{aligned}
\]
Thus, if for each period we develop the function
\SourceLeaf{313}
$X'$ in the way we have shown, we obtain $\nu$ values of each coefficient $\gamma,\gamma',\delta'$ etc., which correspond in succession to the following values of $y$:
\[
\begin{aligned}
&y=\mu a,\\[5pt]
& y=\mu a'=\mu a+\tfrac1\nu360^0,\\[5pt]
& y=\mu a''=\mu a+\tfrac2\nu360^0\text{ etc.}
\end{aligned}
\]
\textit{and therefore form complete periods of $\nu$ terms}. \Needspace{40\baselineskip}
Consequently, by our method, each coefficient $\gamma,\gamma',\delta'$ etc. can be represented by an expression of the form
\[
\begin{aligned}
&\varepsilon+\varepsilon'\cos y+\varepsilon''\cos2y+.\,.\,.\,.+\varepsilon^n\cos ny\\[5pt]
&\phantom{\varepsilon}+\zeta'\sin y+\zeta''\sin2y+.\,.\,.\,.+\zeta^n\sin ny
\end{aligned}
\]
with $n$ either $=\tfrac12\nu-\tfrac12$ or $=\tfrac12\nu$, according as $\nu$ is odd or even. It is evident that such an expression takes the same value on substituting $y=\mu b,\mu c,\mu d$ etc. as it takes for $y=\mu a$, and likewise the same value for $y=\mu b',\mu c',\mu d'$ etc. as for $y=\mu a'$, and so on. Hence, if in
\[
\begin{aligned}
&\gamma+\gamma'\cos x+\gamma''\cos2x+\text{ etc.}\\[5pt]
&\phantom{\gamma}+\delta'\sin x+\delta''\sin2x+\text{ etc.}
\end{aligned}
\]
for the individual coefficients $\gamma,\gamma',\delta'$ etc. these functions of the undetermined angle $y$ are substituted, a function $Z$ involving two angles $x,y$ results, and $Z$ will become the function $X'$ derived from the first period if for $y$ any of the values $\mu a,\mu b,\mu c,\mu d$ etc. is substituted; but if for $y$ any of the values $\mu a',\mu b',\mu c',\mu d'$ etc. is taken, $Z$ will become the function derived from the second period. From this we finally conclude that, if in $Z$ for $y$ we directly write $\mu x$, so that a function of $x$ alone results ($X''$), it will be such that for all $\pi$ values of $x$ it satisfies the proposed values of the function $X$.

\endgroup
